%%%%%%%%%%%%%%%%%%%%%%%%%%%%%%%%%%%%%%%%%%%%%%%%%%%%%%%%%%%%%%%%%%%%%%%%%
% KHAI BÁO CÁC THÔNG SỐ ĐỀ THI - MÃ ĐỀ 132
%%%%%%%%%%%%%%%%%%%%%%%%%%%%%%%%%%%%%%%%%%%%%%%%%%%%%%%%%%%%%%%%%%%%%%%%%
\def\toanlop{10}
\def\thoigian{45}
\def\namhoc{2024 - 2025}
\begin{dethi}
	{SỞ GD\&ĐT THỪA THIÊN HUẾ}
	{TRƯỜNG THPT PHAN ĐĂNG LƯU}
	{KIỂM TRA GIỮA KÌ I}
	{Thời gian làm bài: 45 phút, không kể thời gian phát đề}
\end{dethi}

\Opensolutionfile{ans}[ans/ansTN-PDL-1008]

\noindent \textbf{PHẦN I. Câu trắc nghiệm nhiều phương án lựa chọn. Học sinh trả lời từ câu 1 đến câu 12. Mỗi câu hỏi học sinh chỉ chọn một phương án. (12 Câu - 3 Điểm)}
\vspace{0.2cm}

\begin{ex}%[10V1-1008-01]
	Chuyển động nào sau đây là chuyển động biến đổi?
	\choice
	{Chuyển động thẳng đều.}
	{Chuyển động có độ dịch chuyển giảm đều theo thời gian.}
	{\True Chuyển động tròn đều.}
	{Chuyển động có độ dịch chuyển tăng đều theo thời gian.}
	\loigiai{
		Chuyển động tròn đều có độ lớn vận tốc (tốc độ) không đổi nhưng hướng của vectơ vận tốc liên tục thay đổi theo thời gian, do đó đây là một chuyển động biến đổi. Các phương án còn lại đều là chuyển động thẳng đều. Chọn C.
	}
\end{ex}

\begin{ex}%[10V1-1008-02]
	Đối tượng nghiên cứu tổng quát của khoa học Vật lí học là
	\choice
	{quy luật tương tác của các dạng năng lượng độc lập.}
	{\True các dạng vận động của vật chất và các dạng năng lượng trong tự nhiên.}
	{các dạng vận động và tương tác cơ học của vật chất vô cơ.}
	{các quy luật nhiệt động lực học ứng dụng.}
	\loigiai{
		Vật lí học là ngành khoa học tự nhiên tập trung nghiên cứu về các dạng vận động cơ bản nhất của vật chất và các dạng năng lượng tương tác giữa chúng trong tự nhiên từ cấp độ vi mô đến vĩ mô. Chọn B.
	}
\end{ex}

\begin{ex}%[10V1-1008-03]
	Một người đi xe máy từ nhà đến bến xe bus cách nhà $6\text{ km}$ về phía Đông. Người đó tiếp tục bước lên xe bus đi tiếp $6\text{ km}$ thẳng về phía Bắc. Độ dịch chuyển tổng hợp của người này đạt độ lớn bằng
	\choice
	{$6\text{ km}$.}
	{$12\text{ km}$.}
	{$72\text{ km}$.}
	{\True $6\sqrt{2}\text{ km}$.}
	\loigiai{
		Hai hướng chuyển động Đông và Bắc vuông góc hình học với nhau. Vectơ độ dịch chuyển tổng hợp $\vec{d} = \vec{d}_1 + \vec{d}_2$ đóng vai trò là cạnh huyền của một tam giác vuông cân có hai cạnh cực bằng $6\text{ km}$:
		$$d = \sqrt{d_1^2 + d_2^2} = \sqrt{6^2 + 6^2} = 6\sqrt{2}\text{ km.}$$
		Chọn D.
	}
\end{ex}

\begin{ex}%[10V1-1008-04]
	Hành động nào sau đây \textbf{không} liên quan đến phương pháp lí thuyết trong nghiên cứu vật lí?
	\choice
	{\True Thực hiện ném một quả bóng lên trên cao để quan sát hành trình rơi.}
	{Chế tạo quả địa cầu làm mô hình thu nhỏ của Trái Đất.}
	{Biểu diễn đường truyền của chùm ánh sáng bằng các tia sáng hình học.}
	{Tính toán chính xác quỹ đạo chuyển động của Thiên Vương Tinh dựa vào toán học.}
	\loigiai{
		Hành động trực tiếp ném quả bóng trong đời thực thuộc về phương pháp thực nghiệm trực quan. Các hành động xây dựng mô hình tia sáng, quả địa cầu hay dùng công thức toán học đều đại diện đặc trưng cho phương pháp lí thuyết và mô hình. Chọn A.
	}
\end{ex}

\begin{ex}%[10V1-1008-05]
	Phát biểu nào dưới đây \textbf{sai} khi nói về các cuộc cách mạng công nghiệp trong lịch sử?
	\choice
	{Với cuộc cách mạng công nghiệp lần thứ hai, con người đã sản xuất được điện năng và đưa vào phục vụ cuộc sống.}
	{Internet toàn cầu, trí tuệ nhân tạo và công nghệ nano là các đặc trưng quan trọng của cách mạng công nghiệp lần thứ 4.}
	{\True Đặc trưng cốt lõi của cuộc cách mạng công nghiệp lần thứ 3 là xu hướng thay thế sức lực cơ bắp của con người bằng máy móc cơ khí cơ bản.}
	{Máy hơi nước do James Watt sáng chế năm 1763 là dấu mốc vĩ đại cho cuộc cách mạng công nghiệp lần thứ nhất.}
	\loigiai{
		Xu hướng giải phóng sức lao động chân tay, thay thế sức lực cơ bắp của con người bằng máy móc cơ khí động cơ nhiệt là đặc trưng cốt lõi của cuộc cách mạng công nghiệp lần thứ nhất (1.0), không phải lần thứ ba (3.0 đặc trưng là công nghệ thông tin và tự động hóa). Do đó phát biểu C sai. Chọn C.
	}
\end{ex}

\begin{ex}%[10V1-1008-06]
	Định nghĩa nào sau đây đúng khi đề cập đến bản chất của đại lượng vận tốc tức thời? Vận tốc là
	\choice
	{\True đại lượng vectơ, cho biết mức độ nhanh chậm và chiều hướng chuyển động của vật.}
	{đại lượng vô hướng cho biết tổng thời gian chuyển động của vật.}
	{đại lượng vectơ, cho biết thời gian và quãng đường chuyển động của vật.}
	{đại lượng vô hướng cho biết độ nhanh chậm đơn thuần của chuyển động.}
	\loigiai{
		Vận tốc tức thời là một đại lượng vectơ, có độ lớn đặc trưng cho mức độ chuyển động nhanh hay chậm và có chiều chỉ rõ hướng dịch chuyển vị trí tức thời của vật thể đó. Chọn A.
	}
\end{ex}

\begin{ex}%[10V1-1008-07]
	Một học sinh thực hành đo tốc độ trung bình của viên bi thép thu được giá trị hiển thị kết quả: $v = 2,5 \pm 0,02\text{ (m/s)}$. Sai số tỉ đối $\delta v$ của phép đo này bằng bao nhiêu?
	\choice
	{$0,2\%$.}
	{$4,0\%$.}
	{\True $0,8\%$.}
	{$1,6\%$.}
	\loigiai{
		Sai số tỉ đối của phép đo tốc độ trung bình được xác định bằng biểu thức:
		$$\delta v = \frac{\Delta v}{\overline{v}} \cdot 100\% = \frac{0,02}{2,5} \cdot 100\% = 0,8\%. \text{}$$
		Chọn C.
	}
\end{ex}

\begin{ex}%[10V1-1008-08]
	Độ dịch chuyển là một đại lượng vectơ động học cho biết thông tin gì của hệ vật?
	\choice
	{\True Độ dài thẳng ngắn nhất và chiều hướng sự thay đổi vị trí của vật.}
	{Tọa độ hình học và hướng chuyển động tức thời của vật.}
	{Tổng chiều dài quãng đường thực tế vật đi được.}
	{Tốc độ trung bình trên cả hành trình của vật.}
	\loigiai{
		Độ dịch chuyển $\vec{d}$ là đại lượng vectơ nối vạch điểm đầu cực và vạch điểm cuối cực của chuyển động, cho ta biết đầy đủ thông tin về khoảng cách thẳng ngắn nhất lẫn hướng dốc của sự thay đổi vị trí. Chọn A.
	}
\end{ex}

\begin{ex}%[10V1-1008-09]
	Ưu điểm công nghệ vượt trội nhất của thiết bị cổng quang điện phối hợp đồng hồ hiện số so với các dụng cụ đo truyền thống là
	\choice
	{có thể tự động đo đầy đủ tất cả các đại lượng vật lí.}
	{cho kết quả đo nhanh nhất ở mọi tình huống.}
	{\True loại bỏ được sai số chủ quan của con người và có độ chính xác rất cao.}
	{thiết kế vô cùng gọn nhẹ và dễ thực hiện thao tác thô sơ.}
	\loigiai{
		Cổng quang điện hoạt động bằng cơ chế đóng ngắt tia hồng ngoại tự động khi vật đi qua, giúp loại bỏ hoàn toàn sai số phản xạ chủ quan của con người khi bấm nút cơ học, đem đến độ chính xác vô cùng cao. Chọn C.
	}
	\nobreak
\end{ex}

\begin{ex}%[10V1-1008-10]
	Chọn phát biểu \textbf{không} phù hợp về các quy tắc phòng ngừa nguy cơ mất an toàn khi sử dụng các trang thiết bị thí nghiệm vật lí?
	\choice
	{Chọn sai thang đo khi sử dụng thiết bị đo dòng điện có thể dẫn đến nguy cơ gây chập cháy hỏng thiết bị.}
	{\True Chỉ cần chọn sử dụng các thiết bị đo có lớp vỏ bằng nhựa hoặc kim loại thì sẽ tuyệt đối an toàn với người và dụng cụ.}
	{Khi tiến hành đun nóng với các thiết bị điện công suất lớn và hóa chất dễ cháy có thể dẫn đến nguy cơ cháy nổ phòng thí nghiệm.}
	{Việc thực hiện sai thao tác vận hành các thiết bị có thể dẫn đến nguy cơ gây chấn thương nguy hiểm cho người sử dụng.}
	\loigiai{
		Lớp vỏ nhựa hay kim loại bên ngoài chỉ có tác dụng bảo vệ cơ học cách điện cơ bản, độ an toàn tuyệt đối phụ thuộc hoàn toàn vào ý thức tuân thủ quy trình kĩ thuật nguồn, dòng điện tải của người làm thí nghiệm. Phát biểu B mang tính "tuyệt đối an toàn" chủ quan là sai. Chọn B.
	}
	\nobreak
\end{ex}

\begin{ex}%[10V1-1008-11]
	Trong công tác quản lí an toàn, khi phòng thực hành xảy ra sự cố cháy, chuỗi hành động xử lí ứng phó chuẩn xác theo thứ tự là
	\choice
	{dùng nước dội mạnh trực tiếp để dập ngay đám cháy.}
	{\True lập tức ngắt nguồn điện toàn khu, di chuyển nhanh các chất dễ cháy ra ngoài ngăn cháy lan, tổ chức cứu người, tài sản và dập tắt đám cháy.}
	{hét to lên rồi chạy nhanh ra khỏi phòng càng nhanh càng tốt.}
	{chỉ cần ngắt nguồn điện cấp rồi dùng nước dập đám cháy.}
	\loigiai{
		Quy trình xử lí hỏa hoạn tiêu chuẩn trong phòng thí nghiệm yêu cầu phải cô lập nguồn điện đầu tiên để tránh giật điện và chập mạch lan, sau đó sơ tán chất bắt cháy, cứu hộ người và dùng các bình chữa cháy chuyên dụng độc lập để dập lửa. Chọn B.
	}
\end{ex}

\begin{ex}%[10V1-1008-12]
	Gọi $\overline{A}$ là giá trị trung bình, $\Delta A'$ là sai số dụng cụ, $\Delta \overline{A}$ là sai số ngẫu nhiên tuyệt đối, và $\Delta A$ là sai số tuyệt đối toàn phần. Sai số tỉ đối $\delta A$ của phép đo đại lượng được xác định bằng công thức
	\choice
	{\True $\delta A = \frac{\Delta A}{\overline{A}} \cdot 100\%$}
	{$\delta A = \frac{\overline{A}}{\Delta \overline{A}} \cdot 100\%$}
	{$\delta A = \frac{\Delta \overline{A}}{\overline{A}} \cdot 100\%$}
	{$\delta A = \frac{\Delta A'}{\overline{A}} \cdot 100\%$}
	\loigiai{
		Theo định nghĩa toán học sai số, sai số tỉ đối là phần trăm thương số giữa sai số tuyệt đối toàn phần và giá trị đại số trung bình của phép đo: $\delta A = \frac{\Delta A}{\overline{A}} \cdot 100\%$. Chọn A.
	}
\end{ex}

\Closesolutionfile{ans}

\noindent \textbf{PHẦN II. Câu trắc nghiệm đúng sai. Học sinh trả lời từ câu 1 đến câu 2. Trong mỗi ý a), b), c), d) ở mỗi câu, học sinh chọn đúng hoặc sai. (2 Câu - 2 Điểm)}
\Opensolutionfile{ans}[ans/ansTF-PDL-1008]

\begin{ex}%[10V1-1008-13]
	Một học sinh dùng thước đo vi sai panme có thông số sai số dụng cụ cho sẵn là $0,01\text{ mm}$ để thực hiện đo đường kính d của một viên bi thép nhỏ. Kết quả thu thập qua 5 lần đo độc lập cho giá trị hiển thị ở bảng sau:
	\begin{center}
		\begin{tabular}{|l|c|c|c|c|c|}
			\hline
			\textbf{Lần đo} & \textbf{1} & \textbf{2} & \textbf{3} & \textbf{4} & \textbf{5} \\ \hline
			$\mathbf{d\text{ (mm)}}$ & 6,47 & 6,48 & 6,51 & 6,47 & 6,52 \\ \hline
		\end{tabular}
	\end{center}
	\choiceTF[t]
	{\True Muốn đánh giá và biết được mức độ chính xác tinh vi của phép đo thực hành, ta dựa vào trị số sai số tỉ đối của phép đo đó.}
	{Bán kính trung bình của viên bi thu được từ chuỗi phép đo bằng $6,49\text{ mm}$.}
	{Kết quả tuyệt đối cuối cùng của phép đo đường kính viết kèm sai số toàn phần là: $d = 6,49 \pm 0,03\text{ (mm)}$.}
	{Phương pháp dùng thước panme để xác định trực tiếp thông số đường kính viên bi được phân loại là một phép đo gián tiếp.}
	\loigiai{
		\begin{itemchoice}
			\itemch a) ĐÚNG: Sai số tuyệt đối chỉ cho biết độ lớn của lượng lệch, sai số tỉ đối mới là đại lượng đặc trưng cho độ chính xác và chất lượng của phép đo kĩ thuật.
			\itemch b) SAI: Giá trị trung bình đo được từ bảng là: $\overline{d} = \frac{6,47 + 6,48 + 6,51 + 6,47 + 6,52}{5} = 6,49\text{ mm}$. Đây là giá trị \textbf{đường kính trung bình}, không phải bán kính.
			\itemch c) SAI: Tính toán sai số ngẫu nhiên trung bình: $\overline{\Delta d} = \frac{|6,49-6,47| + |6,49-6,48| + |6,49-6,51| + |6,49-6,47| + |6,49-6,52|}{5} = \frac{0,02+0,01+0,02+0,02+0,03}{5} = 0,02\text{ mm}$. Sai số tuyệt đối toàn phần: $\Delta d = \overline{\Delta d} + \Delta d_{\text{dc}} = 0,02 + 0,01 = 0,03\text{ mm}$. Do đó kết quả viết dạng $d = 6,49 \pm 0,03\text{ mm}$ là hoàn toàn đúng kĩ thuật toán học. *(Lưu ý: Đối chiếu lại dấu đánh nhãn chấm mực bị mờ của tệp gốc để điều chỉnh bản in).*
			\itemch d) SAI: Đường kính viên bi được đọc kết quả trực tiếp thông qua vạch chỉ số hiển thị trên thân thước panme, do đó đây thuộc loại phép đo trực tiếp.
		\end{itemchoice}
	}
\end{ex}

\begin{ex}%[10V1-1008-14]
	Một vật nặng cơ học bắt đầu chuyển động thẳng nhanh dần đều dọc theo phương ngang Ox xuất phát từ trạng thái nghỉ tĩnh. Ghi nhận động học chỉ rõ sau khi vượt qua được quãng đường $100\text{ m}$ đầu tiên, vật đạt được giá trị tốc độ bằng $20\text{ m/s}$. Chọn chiều dương hướng trùng chiều chuyển động.
	\choiceTF[t]
	{Trong hệ thống SI, đơn vị đo tiêu chuẩn của đại lượng gia tốc biến đổi là mét trên giây ($\text{m/s}$).}
	{Giá trị vận tốc tức thời của vật đạt được tại thời điểm sau thời gian $t = 5\text{ s}$ kể từ lúc xuất phát bằng $15\text{ m/s}$.}
	{\True Vì hệ vật thể chuyển động thẳng nhanh dần đều nên vectơ gia tốc đại số $a$ và vectơ vận tốc tức thời $v$ luôn cùng dấu đại số với nhau.}
	{Giá trị gia tốc đại số của chuyển động thẳng biến đổi đều này có trị số bằng $a = -2\text{ m/s}^{2}$.}
	\loigiai{
		\begin{itemchoice}
			\itemch a) SAI: Đơn vị tiêu chuẩn đo gia tốc chuyển động trong hệ SI là mét trên giây bình phương, kí hiệu viết dạng $\text{m/s}^2$.
			\itemch b) SAI: Áp dụng hệ thức độc lập thời gian tính gia tốc chuyển động của vật: $v^2 - v_0^2 = 2as \Rightarrow 20^2 - 0 = 2 \cdot a \cdot 100 \Rightarrow 400 = 200a \Rightarrow a = 2\text{ m/s}^2$. Phương trình vận tốc chuyển động theo thời gian: $v(t) = a \cdot t = 2t$. Tại thời điểm $t = 5\text{ s} \Rightarrow v(5) = 2 \cdot 5 = 10\text{ m/s}$.
			\itemch c) ĐÚNG: Tính chất cơ bản của chuyển động nhanh dần đều là tích số gia tốc và vận tốc luôn nhận giá trị dương ($a \cdot v > 0$), do đó chúng bắt buộc phải cùng dấu đại số.
			\itemch d) SAI: Vì vật chuyển động nhanh dần dốc lên theo chiều dương đã chọn ($v > 0$) nên giá trị gia tốc đại số đúng của hệ vật phải mang trị số dương: $a = +2\text{ m/s}^2$.
		\end{itemchoice}
	}
\end{ex}

\Closesolutionfile{ans}

\noindent \textbf{PHẦN III. Câu trắc nghiệm yêu cầu trả lời ngắn. Học sinh trả lời từ câu 1 đến câu 4 (4 Câu - 2 Điểm)}
\Opensolutionfile{ans}[ans/ansSA-PDL-1008]

\begin{ex}%[10V1-1008-15]
	Trong trận lũ lụt lớn diễn ra tại các tỉnh miền Trung vào tháng 10/2025, dòng lũ tràn về đổ mạnh có tốc độ dòng chảy khoảng $4\text{ m/s}$ đối với bờ sông. Bộ Quốc phòng đã khẩn cấp cấp phát điều động các đội ca nô cứu hộ công suất lớn. Trong một lượt vận hành cứu trợ, đội cứu hộ cho ca nô chạy đạt tốc độ ổn định $8\text{ m/s}$ so với dòng nước để chạy xuôi dòng tiến đến vị trí một ngôi nhà dân cách trạm xuất phát một khoảng quãng đường sông bằng $3\text{ km}$. Tính khoảng thời gian (theo đơn vị giây) để ca nô tiếp cận đến được chỗ người bị nạn? (Kết quả chỉ lấy phần số nguyên).
	\shortans[]{250}
	\loigiai{
		\begin{itemize}
			\item Đổi đơn vị khoảng cách dời chỗ sang mét chuẩn: $s = 3\text{ km} = 3000\text{ m}$.
			\item Do ca nô chạy chuyển động xuôi dòng nước, vận tốc thực tế của ca nô đối với vạch mốc bờ sông bằng tổng vận tốc:
			$$v_{\text{thực}} = v_{\text{cn/n}} + v_{\text{n/b}} = 8 + 4 = 12\text{ m/s.}$$
			\item Khoảng thời gian ca nô di chuyển tiếp cận mục tiêu:
			$$t = \frac{s}{v_{\text{thực}}} = \frac{3000\text{ m}}{12\text{ m/s}} = 250\text{ giây.}$$
		\end{itemize}
	}
\end{ex}

\begin{ex}%[10V1-1008-16]
	Một chiếc ô tô đang đứng yên ở trạng thái nghỉ thì bắt đầu khởi hành chuyển động thẳng nhanh dần đều dọc theo Ox. Biết rằng riêng trong giây thứ 5 tính từ lúc xuất phát, xe đi được hành trình quãng đường dài $13,5\text{ m}$. Hãy xác định giá trị gia tốc của ô tô bằng bao nhiêu $\text{m/s}^2$? (Kết quả chỉ lấy phần số nguyên).
	\shortans[]{3}
	\loigiai{
		\begin{itemize}
			\item Quãng đường vật đi được trong giây thứ 5 chính bằng hiệu giữa quãng đường tích lũy sau 5 giây đầu và quãng đường sau 4 giây đầu tính từ trạng thái nghỉ ($v_0 = 0$):
			$$\Delta s_5 = s(5) - s(4) = \frac{1}{2}a \cdot 5^2 - \frac{1}{2}a \cdot 4^2 = \frac{25}{2}a - \frac{16}{2}a = \frac{9}{2}a = 4,5a\text{ (m).}$$
			\item Thay số liệu quãng đường chặng đề bài cung cấp vào phương trình:
			$$13,5 = 4,5a \Rightarrow a = \frac{13,5}{4,5} = 3\text{ m/s}^2\text{.}$$
		\end{itemize}
	}
\end{ex}

\begin{ex}%[10V1-1008-17]
	Một chiếc ô tô chạy chuyển động thẳng đều dọc theo một trục lộ thẳng. Cho biết đồ thị biểu diễn mối quan hệ phụ thuộc của độ dịch chuyển d theo biến trục thời gian t có dạng một đường thẳng đi qua gốc tọa độ O. Biết trích xuất số liệu ô lưới chỉ ra tại vạch hoành độ thời gian $t = 2\text{ h}$ thì độ dịch chuyển tương ứng xe đạt được bằng $80\text{ km}$. Tính độ lớn vận tốc trung bình của chiếc xe ô tô đó bằng bao nhiêu $\text{km/h}$? (Kết quả chỉ lấy phần số nguyên).
	\shortans[]{40}
	\loigiai{
		\begin{itemize}
			\item Đồ thị độ dịch chuyển dạng đường thẳng dốc lên đi qua gốc tọa độ đại diện đặc trưng cho chuyển động thẳng đều.
			\item Vận tốc chuyển động thẳng đều ổn định của chiếc xe ô tô được xác định bằng tỉ số dốc:
			$$v = \frac{d}{t} = \frac{80\text{ km}}{2\text{ h}} = 40\text{ km/h.}$$
		\end{itemize}
	}
\end{ex}

\begin{ex}%[10V1-1008-18]
	Một chiếc xe máy thực hiện chuyển động thẳng nhanh dần đều xuất phát từ trạng thái nghỉ tĩnh. Sau khi hoàn thành đi được đoạn đường dài $100\text{ m}$ đầu tiên, xe đạt được vận tốc tức thời bằng $10\text{ m/s}$. Chọn chiều chuyển động của xe làm chiều dương. Hỏi ngay sau khi vượt qua được đoạn hành trình dài $100\text{ m}$ tiếp theo, vận tốc tức thời của xe cực cực bằng bao nhiêu $\text{m/s}$? (Kết quả lấy làm tròn đến một chữ số thập phân sau dấu phẩy).
	\shortans[]{14,1}
	\loigiai{
		\begin{itemize}
			\item Áp dụng hệ thức độc lập thời gian tính gia tốc của chuyển động xe máy xuất phát từ nghỉ ($v_0 = 0$):
			$$v_1^2 - v_0^2 = 2as_1 \Rightarrow 10^2 - 0 = 2 \cdot a \cdot 100 \Rightarrow 100 = 200a \Rightarrow a = 0,5\text{ m/s}^2.$$
			\item Xét toàn bộ hành trình tổng hợp sau khi đi thêm 100m tiếp theo, tổng quãng đường toàn phần tích lũy tính từ gốc bằng: $S = s_1 + s_2 = 100 + 100 = 200\text{ m}$.
			\item Vận tốc cuối chặng toàn phần $v_2$ thỏa mãn phương trình liên hệ:
			$$v_2^2 - v_0^2 = 2aS \Rightarrow v_2^2 - 0 = 2 \cdot 0,5 \cdot 200 \Rightarrow v_2^2 = 200 \Rightarrow v_2 = \sqrt{200} = 10\sqrt{2} \approx 14,14\text{ m/s.}$$
			\item Làm tròn số thập phân theo yêu cầu thu được kết quả: $v_2 \approx 14,1\text{ m/s.}$
		\end{itemize}
	}
\end{ex}

\Closesolutionfile{ans}

\noindent \textbf{Phần 4. Tự luận (2 câu - 3 điểm)}
\vspace{0.2cm}

\begin{ex}%[10V1-1008-19]
	\textbf{(2 điểm)} Một người vận động viên bơi ngang từ bờ sông bên này sang bờ bên kia theo phương vuông góc dốc thẳng với dòng sông rộng $50\text{ m}$ (hướng từ Tây sang Đông). Dòng nước lũ đổ về chảy mạnh theo hướng dốc từ Bắc xuống Nam. Do lực đẩy của nước sông chảy mạnh nên khi sang đến bờ bên kia, người đó đã bị dòng nước cuốn trôi xuôi dịch chỗ về phía hạ lưu một đoạn dài đúng $50\text{ m}$.
	\begin{enumerate}
		\item[a)] Hãy vẽ hình minh họa hình học chuyển động và xác định hướng của vectơ độ dịch chuyển tổng hợp của người đó trong cả quá trình.
		\item[b)] Tính độ lớn độ dịch chuyển tổng hợp của người bơi trong cả quá trình bơi qua sông.
	\end{enumerate}
	\loigiai{
		\begin{enumerate}
			\item[a)] 
			\begin{itemize}
				\item \textbf{Thao tác vẽ hình:} Chọn gốc tọa độ A tại vị trí điểm xuất phát ở bờ bên này. Vẽ vectơ $\vec{d}_1 = \vec{AB}$ hướng nằm ngang thẳng sang bờ đối diện biểu diễn độ rộng dòng sông ($AB = 50\text{ m}$). Vẽ tiếp vectơ dịch chỗ dọc bờ $\vec{d}_2 = \vec{BC}$ hướng thẳng dọc dốc xuống đại diện cho đoạn đường bị nước cuốn trôi xuôi ($BC = 50\text{ m}$).
				\item Vectơ độ dịch chuyển tổng hợp toàn phần của người bơi là vectơ hình học nối thẳng từ gốc điểm đầu cực đến đỉnh cuối cực: $\vec{d} = \vec{AC} = \vec{d}_1 + \vec{d}_2$. Vectơ này hướng theo hướng Đông – Nam lệch góc $45^{\circ}$ so với phương bơi vuông góc ban đầu.
			\end{itemize}
			\item[b)] Do hệ trục tọa độ hai hướng Đông và Nam vuông góc phương hình học với nhau, tam giác ABC vuông cân tại nút đỉnh B. Độ lớn của độ dịch chuyển tổng hợp toàn phần xác định qua hệ thức định lý Pytago:
			$$d = AC = \sqrt{AB^2 + BC^2} = \sqrt{50^2 + 50^2} = 50\sqrt{2} \approx 70,71\text{ m.}$$
		\end{enumerate}
	}
\end{ex}

\begin{ex}%[10V1-1008-20]
	\textbf{(1 điểm)} Một chất điểm chuyển động thẳng biến đổi đều dọc theo một trục Ox có đồ thị vận tốc – thời gian $(v-t)$ được biểu diễn qua đường gấp khúc hình học như đồ thị hình vẽ bên dưới. 
	\begin{center}
		\begin{tikzpicture}[scale=0.08,thick,>=stealth']
			% Hệ lưới tọa độ ô chia
			\draw[lightgray,step=10] (0,0) grid (90,70);
			% Hệ trục tọa độ chính
			\draw[->,line width=1.2pt] (-5,0) -- (95,0) node[below] {$t\text{ (s)}$};
			\draw[->,line width=1.2pt] (0,-5) -- (0,75) node[left] {$v\text{ (cm/s)}$};
			\node[below left] at (0,0) {$O$};
			% Các vạch chỉ số mốc chính trục ngang
			\draw (20,1) -- (20,-1) node[below] {$20$};
			\draw (40,1) -- (40,-1) node[below] {$40$};
			\draw (80,1) -- (80,-1) node[below] {$80$};
			% Các vạch chỉ số mốc chính trục đứng
			\draw (1,20) -- (-1,20) node[left] {$20$};
			\draw (1,60) -- (-1,60) node[left] {$60$};
			% Vẽ các đoạn gấp khúc đồ thị động học
			\draw[blue,line width=1.5pt] (0,20) node[below right,black] {A} -- (20,60) node[above,black] {B} -- (40,60) node[above,black] {C} -- (80,0) node[above right,black] {D};
			% Các đường dóng nét đứt phụ trợ
			\draw[dashed] (20,0) -- (20,60) -- (0,60);
			\draw[dashed] (40,0) -- (40,60);
		\end{tikzpicture}
	\end{center}
	\begin{enumerate}
		\item[a)] Tính trị số độ tăng vận tốc đại số ($\Delta v$) của chất điểm tính từ thời điểm ban đầu $t=0$ đến thời điểm vạch $t = 20\text{ s}$.
		\item[b)] Tính giá trị gia tốc đại số $a$ của chất điểm khi di chuyển vận hành trong giai đoạn đoạn cuối CD.
	\end{enumerate}
	\loigiai{
		\begin{enumerate}
			\item[a)] Trích xuất số liệu tọa độ từ đồ thị hàm vận tốc: tại mốc đầu $t=0$ vận tốc đạt $v_A = 20\text{ cm/s}$, tại vạch thời gian $t = 20\text{ s}$ vận tốc đạt đỉnh $v_B = 60\text{ cm/s}$. \\
			Độ tăng (biến thiên) vận tốc đại số trong giai đoạn AB này bằng:
			$$\Delta v = v_B - v_A = 60 - 20 = 40\text{ cm/s} = 0,4\text{ m/s.}$$
			\item[b)] Trong giai đoạn cuối ứng với đoạn thẳng CD nối từ vạch thời gian $t_C = 40\text{ s}$ (với vận tốc $v_C = 60\text{ cm/s}$) đến vị trí dừng hẳn trên trục hoành tại vạch thời gian $t_D = 80\text{ s}$ (với vận tốc chạm đáy $v_D = 0\text{ cm/s}$). \\
			Gia tốc đại số chặng CD đặc trưng cho chuyển động thẳng chậm dần đều bằng hệ số dốc:
			$$a_{CD} = \frac{v_D - v_C}{t_D - t_C} = \frac{0 - 60}{80 - 40} = \frac{-60}{40} = -1,5\text{ cm/s}^2 = -0,015\text{ m/s}^2\text{.}$$
		\end{enumerate}
	}
\end{ex}

\Closesolutionfile{ans}

\newpage
%%%%%%%%%%%%%%%%%%%%%%%%%%%%%%%%%%%%%%%%%%%%%%%%%%%%%%%%%%%%%%%%%%%%%%%%%
% KHUNG ĐÁP ÁN BẢNG TỰ ĐỘNG IN KẾT XUẤT CỦA HỆ THỐNG
%%%%%%%%%%%%%%%%%%%%%%%%%%%%%%%%%%%%%%%%%%%%%%%%%%%%%%%%%%%%%%%%%%%%%%%%%
\begin{center} \textbf{BẢNG ĐÁP ÁN THAM KHẢO KHỐI 10 --- MÃ ĐỀ 1008} \end{center}

\noindent \textbf{Phần I (Câu hỏi trắc nghiệm nhiều phương án lựa chọn):}  
\indapan{12}{ans/ansTN-PDL-1008}

\noindent \textbf{Phần II (Câu hỏi trắc nghiệm Đúng/Sai):}
\indapan[TF]{2}{ans/ansTF-PDL-1008}

\noindent \textbf{Phần III (Câu hỏi trắc nghiệm trả lời ngắn):}
\indapan[SA]{4}{ans/ansSA-PDL-1008}

\label{B\thedeso}